\documentclass[11pt]{amsart}
\usepackage[T1]{fontenc}
\usepackage[utf8]{inputenc}
\usepackage{amsmath,amssymb,amsthm}
\usepackage{mathtools}
\usepackage[margin=1.1in]{geometry}

\providecommand{\C}{\mathbb{C}}
\providecommand{\bC}{\mathbb{C}}
\providecommand{\R}{\mathbb{R}}
\providecommand{\Z}{\mathbb{Z}}
\providecommand{\dbar}{\bar{\partial}}
\providecommand{\g}{\mathfrak{g}}
\providecommand{\Obs}{\mathrm{Obs}}
\providecommand{\Sym}{\mathrm{Sym}}
\providecommand{\gh}{\mathrm{gh}}
\providecommand{\af}{\mathrm{af}}
\providecommand{\kappach}{\kappa_{\mathrm{ch}}}

\theoremstyle{plain}
\newtheorem{theorem}{Theorem}[section]
\newtheorem{proposition}[theorem]{Proposition}
\newtheorem{lemma}[theorem]{Lemma}
\theoremstyle{definition}
\newtheorem{definition}[theorem]{Definition}
\newtheorem{construction}[theorem]{Construction}
\theoremstyle{remark}
\newtheorem{remark}[theorem]{Remark}

\title{Abelian holomorphic Chern--Simons on $\C^3$:\\
the BV--BRST complex, propagator, defect VOA}
\author{Wave-12 A5/20}
\date{2026-04-22}

\begin{document}
\maketitle

Let $X = \C^3$ with coordinates $(z_1, z_2, z_3)$ and holomorphic volume form
$\Omega = dz_1 \wedge dz_2 \wedge dz_3$. Write $\mathcal{E} = \Omega^{0,\bullet}(X)$ as
a graded vector space, the Dolbeault complex of the trivial line bundle.

\section{The BV--BRST complex}

\begin{construction}[Field content and degree data]
\label{cons:bv-field}
The BV superfield for Abelian hCS is
\[
\mathcal{A} \;=\; c \;+\; A_{0,1} \;+\; A^{\ast}_{0,2} \;+\; c^{\ast}_{0,3}
\;\in\; \Omega^{0,\bullet}(\C^3).
\]
The assignment of ghost number $\gh$, antifield number $\af$, and total BV
degree $\deg_{\mathrm{BV}} = \gh - \af$ is
\[
\begin{array}{c|cccc}
 & c & A_{0,1} & A^{\ast}_{0,2} & c^{\ast}_{0,3} \\ \hline
\text{Dolbeault } (0,q) & 0 & 1 & 2 & 3 \\
\gh & +1 & 0 & -1 & -2 \\
\af & 0 & 0 & 1 & 2 \\
\deg_{\mathrm{BV}} & +1 & 0 & -2 & -4
\end{array}
\]
\end{construction}

The total BV grading is $\deg_{\mathrm{BV}} = (\text{Dolbeault degree}) - 1 - 2\af$,
equivalently $\gh$ is $1$ minus Dolbeault degree on the physical
pair $(c, A_{0,1})$ and the antifield pair is obtained by Koszul duality,
shifted to satisfy the $(-1)$-shifted symplectic pairing of
Costello--Gwilliam.

\begin{construction}[$(-1)$-shifted symplectic pairing]
\label{cons:bv-pairing}
The BV antibracket arises from the non-degenerate pairing of cohomological
degree $-1$
\[
\omega_{\mathrm{BV}}(\mathcal{A}_1, \mathcal{A}_2) \;=\;
\int_{\C^3} \Omega \wedge \mathcal{A}_1 \wedge \mathcal{A}_2,
\]
which pairs $(c, c^{\ast}_{0,3})$ and $(A_{0,1}, A^{\ast}_{0,2})$ off-diagonally
by Serre duality on $\Omega^{0,\bullet}$.
\end{construction}

\section{Abelian action and BRST differential}

\begin{proposition}[Quadratic action]
\label{prop:quadratic-action}
For $G = U(1)$ the cubic term $\tfrac{2}{3} \mathcal{A} [\mathcal{A}, \mathcal{A}]$
vanishes and
\[
S_{\mathrm{hCS}}(\mathcal{A}) \;=\; \int_{\C^3} \Omega \wedge \mathcal{A} \,\dbar\, \mathcal{A}.
\]
\end{proposition}

\begin{proof}
Abelian bracket $[\cdot,\cdot] = 0$. The residual kinetic term, restricted to
the physical pair, recovers the Dolbeault action
$\int \Omega \wedge A_{0,1} \dbar A_{0,1}$ of Costello 2013, §2
(arXiv:1303.2632), with ghost $c$ implementing the gauge symmetry
$A_{0,1} \mapsto A_{0,1} + \dbar c$.
\end{proof}

\begin{definition}[BRST differential $Q$]
\label{def:Q-BRST}
On generators, $Q$ is the classical BV differential associated to
$S_{\mathrm{hCS}}$,
\[
Q c \;=\; 0, \qquad
Q A_{0,1} \;=\; \dbar c, \qquad
Q A^{\ast}_{0,2} \;=\; \dbar A_{0,1}, \qquad
Q c^{\ast}_{0,3} \;=\; \dbar A^{\ast}_{0,2}.
\]
Equivalently $Q = \dbar$ on the total complex
$(\Omega^{0,\bullet}(\C^3), \dbar)$ with the degree re-labelling of
Construction~\ref{cons:bv-field}.
\end{definition}

$Q^2 = \dbar^2 = 0$. The Abelian BV complex is a shifted copy of the
Dolbeault complex of $\C^3$; its cohomology is the polynomial ring of
holomorphic functions placed in ghost number $+1$ (the $c$-sector).

\section{Propagator}

\begin{proposition}[Holomorphic propagator]
\label{prop:propagator}
The inverse of $\dbar$ on $\Omega^{0,\bullet}(\C^3)$, modulo harmonic
representatives, is given by the Bochner--Martinelli kernel
\[
P(z, w) \;=\; \frac{(n-1)!}{(2\pi i)^{n}}
\sum_{k=1}^{n}
\frac{(-1)^{k-1} (\overline{z_k - w_k})\,
d\bar{z}_1 \wedge \widehat{d\bar{z}_k} \wedge \cdots \wedge d\bar{z}_n}
{|z - w|^{2n}}
\]
with $n = 3$, so that $\dbar_w P(z, w) = \delta_{z=w}$ as a current of
Dolbeault bidegree $(n, n-1) = (3, 2)$ on the diagonal-complement.
\end{proposition}

\begin{proof}
Bochner--Martinelli on $\C^n$; see Costello--Gwilliam 2017 Vol II §5.2.
Contraction of $P$ against $\Omega$ produces the scalar two-point function
$\langle A_{0,1}(z) A^{\ast}_{0,2}(w) \rangle$ governing perturbative
hCS.
\end{proof}

\section{Observables}

\begin{construction}[Classical observables at leading order]
\label{cons:obs-classical}
The classical observable algebra is
\[
\Obs^{\mathrm{cl}}(\C^3) \;=\; \Sym\bigl( \mathcal{E}^{\vee}[-1] \bigr),
\]
the free graded-commutative algebra on linear functionals of the
BV fields, with differential induced by $\dbar$. The $H^0$ classes are
Dolbeault-holomorphic functions of $A_{0,1}$; perturbative quantisation
(Costello--Gwilliam Vol II, Theorem 5.4) deforms $\Obs^{\mathrm{cl}}$ to
$\Obs^{\mathrm{q}}$ via the propagator of
Proposition~\ref{prop:propagator}. For Abelian $G$ the quantum
correction is Gaussian: $\Obs^{\mathrm{q}}$ is a free holomorphic boson
on $\C^3$, i.e.\ the $E_1$-algebra of polynomial functionals of
$A_{0,1}$ with Wick contractions governed by $P(z,w)$.
\end{construction}

\section{Line defect and defect VOA}

\begin{construction}[Line defect on $\C_{z_1} \subset \C^3$]
\label{cons:line-defect}
Let $L = \C_{z_1} = \{z_2 = z_3 = 0\}$. A line defect is a morphism
$\phi \colon \Obs^{\mathrm{q}}(\C^3 \setminus L) \to \mathcal{V}$ to a
holomorphic factorisation algebra on $\C_{z_1}$. Restricting the
Bochner--Martinelli propagator transversely to $L$ yields a $2$-dimensional
holomorphic propagator in the normal $(z_2, z_3)$ directions, reducing the
bulk field to a $\beta\gamma$-system on $L$:
\[
\beta(z_1) \;:=\; A_{0,1}|_{L}, \qquad \gamma(z_1) \;:=\; c|_{L},
\qquad \langle \beta(z) \gamma(w) \rangle \;=\; \frac{1}{z-w}.
\]
The defect VOA is
\[
\mathcal{V}_L \;=\; \beta\gamma\text{-system on } \C_{z_1},
\]
with central charge $c = 2$. This recovers the Costello--Li 2016
identification (arXiv:1606.00365, Theorem 1.1) of the Abelian hCS
line-defect algebra as the rank-one symplectic boson.
\end{construction}

\begin{remark}[Non-Abelian extension]
\label{rem:non-ab}
For $G$ non-Abelian the cubic vertex $\tfrac{2}{3} \mathcal{A} [\mathcal{A}, \mathcal{A}]$
reinstates and the defect VOA is the affine Kac--Moody vertex algebra
at level determined by the one-loop anomaly; for
$G = GL_n$ one recovers $V^{\mathrm{crit}}(\widehat{\g})$ (Costello 2013,
Theorem 9.2.1). The Abelian case is the $n = 1$ free-field
specialisation.
\end{remark}

\begin{remark}[$\kappach$-scope]
\label{rem:kappa-scope}
The BV pairing of Construction~\ref{cons:bv-pairing} computes
$\kappach(\C^3) = \chi(\mathcal{O}_{\C^3}) = 1$ as the constant term of
the Hodge supertrace, consistent with the $d=3$ stratification of
\texttt{chapters/examples/cy\_d\_kappa\_stratification.tex}. The line
defect selects a $d=1$ child whose supertrace matches the $\beta\gamma$
central charge $c = 2 \equiv 2\kappach$.
\end{remark}

\paragraph{Provenance.}
Costello 2013 (arXiv:1303.2632) §2--3: hCS on $\C^n$, BV formulation.
Costello--Gwilliam 2017 Vol II §5.1--5.4: Abelian holomorphic
factorisation algebras, Bochner--Martinelli propagator, free
$\beta\gamma$-system.
Costello--Li 2016 (arXiv:1606.00365) Theorems 1.1 and 3.2: defect VOA
on line in $\C^3$.

\end{document}
